\documentclass[11pt,a4paper]{article}
\usepackage[T1]{fontenc}
\usepackage[utf8]{inputenc}
\usepackage{amsmath,amssymb,xcolor}
\usepackage[margin=25mm]{geometry}
\usepackage[colorlinks=true,urlcolor=blue]{hyperref}
\definecolor{nexorablue}{HTML}{2459B5}
\setlength{\parindent}{0pt}
\setlength{\parskip}{7pt}
\setlength{\emergencystretch}{2em}
\renewcommand{\familydefault}{\sfdefault}
\begin{document}

{\small\bfseries\color{nexorablue} NEXORA CPGE}\par
{\small MPSI / NEX-CIR-01 / Corrigé}\par
{\LARGE\bfseries Propagation et courants signés}\par
\vspace{0.25em}\hrule\vspace{0.6em}

\section*{Ordres de grandeur}
Le temps de propagation vaut $0{,}75\,\mathrm{ns}$. Pour $2\,\mathrm{MHz}$, la période vaut $500\,\mathrm{ns}$ et le rapport vaut $0{,}0015$. Pour $500\,\mathrm{MHz}$, la période vaut $2\,\mathrm{ns}$ et le rapport vaut $0{,}375$. L’ARQS est bien justifiée par ce critère dans le premier cas. Pour le second, il faut tenir compte du retard de propagation.

\section*{Bilan de charge}
Les valeurs en milliampères vérifient $3{,}5-1{,}2=0{,}8+i_4$, donc $i_4=1{,}5\,\mathrm{mA}$. Le courant associé à $-1{,}2\,\mathrm{mA}$ quitte réellement le nœud, bien que sa flèche soit entrante.


\end{document}
